\documentclass[11pt,a4paper]{article}

% --- Mathematics (must precede unicode-math) ---
\usepackage{amsmath}
\usepackage{mathtools}

% --- Fonts (lualatex) ---
\usepackage{fontspec}
\usepackage{unicode-math}
\setmathfont{KpMath-Regular.otf}
\setmathfont{KpMath-Bold.otf}[version=bold]

% --- Microtypography (after all font setup) ---
\usepackage{microtype}

% --- Colors & page layout ---
\usepackage[dvipsnames]{xcolor}
\usepackage[margin=25mm, top=30mm, bottom=30mm]{geometry}

% --- Headers & footers ---
\usepackage{fancyhdr}
\pagestyle{fancy}
\fancyhf{}
\fancyhead[L]{\small\itshape Cubic T-Matrix: Full Green's Tensor}
\fancyhead[R]{\small\itshape Nestor (1996)}
\fancyfoot[C]{\thepage}
\renewcommand{\headrulewidth}{0.4pt}
\renewcommand{\footrulewidth}{0pt}

% --- Section numbering ---
\setcounter{secnumdepth}{3}

% --- Tables ---
\usepackage{booktabs}
\usepackage{array}

% --- Captions ---
\usepackage[font=small, labelfont=bf, labelsep=period]{caption}

% --- Spacing ---
\usepackage{setspace}
\setstretch{1.15}
\setlength{\parskip}{0.4em}

% --- Citations ---
\usepackage[round,authoryear]{natbib}
\bibliographystyle{plainnat}

% --- Hyperlinks (last before macros) ---
\usepackage[colorlinks=true, linkcolor=blue!60!black,
            citecolor=blue!60!black, urlcolor=blue!60!black]{hyperref}

% --- Macros (reused from nearfield_verification.tex) ---
\newcommand{\bG}{\mathbf{G}}
\newcommand{\bI}{\mathbf{I}}
\newcommand{\bx}{\mathbf{x}}
\newcommand{\beR}{\hat{\mathbf{e}}_R}
\newcommand{\dd}{\mathrm{d}}
\newcommand{\order}[1]{\mathcal{O}\!\left(#1\right)}

% --- New macros for this document ---
\newcommand{\bT}{\mathbf{T}}
\newcommand{\bE}{\mathbf{E}}
\newcommand{\bS}{\mathbf{S}}
\newcommand{\Dlam}{\Delta\lambda}
\newcommand{\Dmu}{\Delta\mu}
\newcommand{\Drho}{\Delta\rho}
\newcommand{\PV}{\operatorname{P.V.}}
\newcommand{\re}{\operatorname{Re}}
\newcommand{\im}{\operatorname{Im}}

\begin{document}

% --- Title block ---
\thispagestyle{plain}
\begin{center}
\vspace*{1em}
{\LARGE\bfseries Cubic T-Matrix with the Full\\[0.3em]
Elastodynamic Green's Tensor}\\[1.8em]
{\large T.\,M.\ Nestor}\\[0.5em]
{\itshape Research School of Earth Sciences, Australian National University}\\[2em]
\end{center}

% =============================================================================
\section{Introduction}
% =============================================================================

The T-matrix formulation of elastic wave scattering expresses the
scattered field from an inclusion as a product of a scattering
operator (the T-matrix) and the incident field
\citep{GubernatisDomanyKrumhansl1977}. For inclusions whose
dimensions are small compared to the wavelength ($ka \ll 1$, the
Rayleigh regime), the T-matrix depends on the volume integral of
the elastodynamic Green's tensor over the scatterer
\citep{Nestor1996}.

For spherical inclusions, this integral inherits full rotational
symmetry: the depolarization tensor has only two independent
components, and the T-matrix decomposes into volumetric and shear
parts. For a \emph{cube}, however, the cubic point-group symmetry
($O_h$) is lower than spherical, and a third independent component
appears---the \emph{cubic anisotropy} term. This term is governed
by the fourth-rank tensor
\begin{equation}
  E_{ijkl} = \sum_{m=1}^{3}
    \delta_{im}\,\delta_{jm}\,\delta_{km}\,\delta_{lm},
  \label{eq:Etensor}
\end{equation}
which is nonzero only when all four indices are equal.

The complete derivation requires:
\begin{enumerate}
  \item The \emph{static} Eshelby depolarization tensor for the
    cube \citep{Eshelby1957}, obtained from the Kelvin (static)
    Green's tensor via surface integrals involving exact geometric
    constants.
  \item The \emph{dynamic} radiation corrections, obtained from the
    smooth polynomial part of the full elastodynamic Green's tensor.
\end{enumerate}
Previous treatments \citep{Nestor1996,GubernatisDomanyKrumhansl1977b}
often presented only the static (Born) limit. This document
presents the complete analytic result, with the full
frequency-dependent Green's tensor, verified symbolically in
Mathematica (see accompanying notebook
\texttt{CubicTMatrix\_FullGreensTensor.nb}).

% =============================================================================
\section{The Elastodynamic Green's Tensor}
\label{sec:greens}
% =============================================================================

\subsection{Kupradze representation}

The frequency-domain Green's tensor satisfies
\begin{equation}
  \mu\,\nabla^2 G_{ij}
  + (\lambda + \mu)\,\partial_i\partial_k G_{kj}
  + \rho\omega^2 G_{ij}
  = -\delta_{ij}\,\delta(\bx),
  \label{eq:pde}
\end{equation}
with the Kupradze representation \citep{Kupradze1963,Kausel2006}
detailed in the companion document \citep[see][for
derivation]{WuBenMenahem1985}.

\subsection{Radial decomposition}

For an isotropic homogeneous medium, the Green's tensor admits the
radial decomposition
\begin{equation}
  \boxed{
  G_{ij}(\bx) = f(r)\,\delta_{ij}
    + \frac{g(r)}{r^2}\,x_i\,x_j,
  }
  \label{eq:radial}
\end{equation}
where $r = |\bx|$, and $f(r)$ and $g(r)$ encode the P-wave, S-wave,
and near-field contributions.

Define the near-field constant
\begin{equation}
  C_\text{nf}
  = \frac{1 - \beta^2/\alpha^2}{8\pi\rho\beta^2}
  = \frac{\lambda + \mu}{8\pi\mu(\lambda + 2\mu)},
  \label{eq:Cnf}
\end{equation}
and the P-wave and S-wave radial functions
\begin{equation}
  X(r) = \frac{e^{i\omega r/\alpha}}{4\pi\rho\alpha^2},
  \qquad
  V(r) = \frac{e^{i\omega r/\beta}}{4\pi\rho\beta^2}.
  \label{eq:XV}
\end{equation}
Then
\begin{equation}
  f(r) = \frac{V(r) - C_\text{nf}}{r},
  \qquad
  g(r) = \frac{3\,C_\text{nf} + X(r) - V(r)}{r}.
  \label{eq:fg}
\end{equation}

The static ($\omega \to 0$) limit gives the Kelvin Green's tensor:
\begin{equation}
  G^0_{ij}(\bx) = \frac{1}{r}
  \biggl[a_0\,\delta_{ij} + b_0\,\frac{x_i x_j}{r^2}\biggr],
  \label{eq:kelvin}
\end{equation}
with
\begin{equation}
  \boxed{
  a_0 = \frac{\alpha^2 + \beta^2}{8\pi\rho\alpha^2\beta^2},
  \qquad
  b_0 = \frac{\alpha^2 - \beta^2}{8\pi\rho\alpha^2\beta^2}.
  }
  \label{eq:a0b0}
\end{equation}

% =============================================================================
\section{Smooth Polynomial Decomposition}
\label{sec:smooth}
% =============================================================================

The key insight is that $f(r) \cdot r$ and $g(r) \cdot r$ are
entire functions of~$r$, expandable in Taylor series. When
decomposed by parity:
\begin{itemize}
  \item \textbf{Even powers} of $r$ (constant, $r^2$, $r^4$,
    \ldots) contribute the static singular ($1/r$) part after
    dividing by $r$---these yield the Eshelby depolarization.
  \item \textbf{Odd powers} of $r$ ($r$, $r^3$, $r^5$,
    \ldots) remain smooth polynomials after dividing by $r$---these
    yield the radiation corrections.
\end{itemize}

\subsection{Taylor expansion}

Expanding the exponentials to order~$N$:
\begin{align}
  f(r) \cdot r &= V(r) - C_\text{nf}
  = \sum_{n=0}^{N} c_n^{(f)}\, r^n, \label{eq:fser} \\
  g(r) \cdot r &= 3C_\text{nf} + X(r) - V(r)
  = \sum_{n=0}^{N} c_n^{(g)}\, r^n. \label{eq:gser}
\end{align}
The even-order coefficients reproduce the static Green's tensor;
the odd-order coefficients give the smooth (purely imaginary at
leading order) corrections.

\subsection{Smooth Green's tensor}

The smooth part of the Green's tensor takes the form
\begin{equation}
  G^s_{ij}(\bx) = \Phi(r^2)\,\delta_{ij}
    + \Psi(r^2)\,x_i\,x_j,
  \label{eq:Gsmooth}
\end{equation}
where
\begin{equation}
  \Phi(u) = \sum_{n=0}^{N_\Phi} \Phi_n\, u^n,
  \qquad
  \Psi(u) = \sum_{n=0}^{N_\Psi} \Psi_n\, u^n,
  \label{eq:PhiPsi}
\end{equation}
with $u = r^2 = x_1^2 + x_2^2 + x_3^2$.

The leading coefficients, verified symbolically in Mathematica,
are:

\begin{table}[ht]
\centering
\caption{Leading Taylor coefficients $\Phi_n$ and $\Psi_n$ of the
  smooth Green's tensor.}
\label{tab:coeffs}
\begin{tabular}{@{}cll@{}}
\toprule
$n$ & $\Phi_n$ & $\Psi_n$ \\
\midrule
0 & $\displaystyle\frac{i\omega}{4\pi\beta^3\rho}$
  & $\displaystyle\frac{i(\alpha^5 - \beta^5)\omega^3}
     {24\pi\alpha^5\beta^5\rho}$ \\[12pt]
1 & $\displaystyle\frac{-i\omega^3}{24\pi\beta^5\rho}$
  & $\displaystyle\frac{-i(\alpha^7 - \beta^7)\omega^5}
     {480\pi\alpha^7\beta^7\rho}$ \\[12pt]
2 & $\displaystyle\frac{i\omega^5}{480\pi\beta^7\rho}$
  & $\displaystyle\frac{i(\alpha^9 - \beta^9)\omega^7}
     {20160\pi\alpha^9\beta^9\rho}$ \\[12pt]
3 & $\displaystyle\frac{-i\omega^7}{20160\pi\beta^9\rho}$
  & $\displaystyle\frac{-i(\alpha^{11} - \beta^{11})\omega^9}
     {1451520\pi\alpha^{11}\beta^{11}\rho}$ \\
\bottomrule
\end{tabular}
\end{table}

\noindent
All $\Phi_n$ and $\Psi_n$ are purely imaginary, reflecting the
radiative (energy-losing) character of the smooth contribution.

% =============================================================================
\section{Volume Integrals over the Cube}
\label{sec:cube}
% =============================================================================

Consider a cube of half-side~$a$ centred at the origin:
$V = [-a,a]^3$ with volume $|V| = 8a^3$.

\subsection{Monomial moments}

By symmetry, all odd-order moments vanish. The nonzero even moments
relevant to the T-matrix are:
\begin{equation}
  \int_V x_i^2 \,\dd V = \frac{8a^5}{3},
  \qquad
  \int_V x_i^4 \,\dd V = \frac{8a^7}{5},
  \qquad
  \int_V x_i^2 x_j^2 \,\dd V = \frac{8a^7}{9}
  \quad (i \neq j).
  \label{eq:moments}
\end{equation}

\subsection{Fourth-rank moment tensor}

The fourth-order moment
$M_{ijkl} = \int_V x_i\,x_j\,x_k\,x_l\,\dd V$
decomposes under cubic symmetry into isotropic and anisotropic
parts:
\begin{equation}
  M_{ijkl} = P\,(\delta_{ij}\delta_{kl}
    + \delta_{ik}\delta_{jl} + \delta_{il}\delta_{jk})
  + Q\,E_{ijkl},
  \label{eq:moment4}
\end{equation}
where
\begin{equation}
  P = \frac{8a^7}{9},
  \qquad
  Q = 8a^7\!\left(\frac{1}{5} - \frac{3}{9}\right)
    = -\frac{16a^7}{45}.
  \label{eq:PQ}
\end{equation}
The \emph{isotropy ratio}
$M_{iiii}/(3\,M_{iijj}) = 3/5$ (compared to $1$ for a sphere)
quantifies the departure from spherical symmetry.

\subsection{The cubic anisotropy tensor}

The tensor~$E_{ijkl}$ defined in~\eqref{eq:Etensor} satisfies:
\begin{itemize}
  \item Full symmetry: $E_{ijkl}$ is invariant under any
    permutation of indices.
  \item Trace: $E_{iikl} = \delta_{kl}$.
  \item Contraction with isotropic tensors:
    $E_{ijkl}\,\delta_{kl} = \delta_{ij}$.
\end{itemize}
It is the unique fourth-rank tensor (up to scale) that distinguishes
cubic from spherical symmetry in the decomposition of
$I_{ijkl}$.

% =============================================================================
\section{The Integral Tensor \texorpdfstring{$I_{ijkl}$}{I\_ijkl}}
\label{sec:integral}
% =============================================================================

The central quantity for the T-matrix is the volume integral of the
symmetrised second derivatives of the Green's tensor:
\begin{equation}
  I_{ijkl} = \PV \int_V
    \frac{\partial^2 G_{ij}(\bx)}{\partial x_k \,\partial x_l}
    \,\dd V,
  \label{eq:Idef}
\end{equation}
where the principal value excludes the origin singularity. By
cubic symmetry, this decomposes as
\begin{equation}
  \boxed{
  I_{ijkl} = A\,\delta_{ij}\delta_{kl}
    + B\,(\delta_{ik}\delta_{jl} + \delta_{il}\delta_{jk})
    + C\,E_{ijkl}.
  }
  \label{eq:Idecomp}
\end{equation}
For a sphere, $C = 0$ and only two independent constants remain.
For the cube, $C \neq 0$, and this is the origin of cubic anisotropy
in the T-matrix.

The three coefficients are determined from the independent
components:
\begin{equation}
  A = I_{1122}, \qquad
  B = I_{1212}, \qquad
  C = I_{1111} - I_{1122} - 2\,I_{1212}.
  \label{eq:ABC_from_I}
\end{equation}

Each coefficient receives a static (Eshelby) and smooth (radiation)
contribution:
\begin{equation}
  A = A_\text{stat} + A_\text{smooth}, \qquad
  B = B_\text{stat} + B_\text{smooth}, \qquad
  C = C_\text{stat} + C_\text{smooth}.
  \label{eq:split}
\end{equation}

% -----------------------------------------------------------------------------
\subsection{Static Eshelby depolarization tensor}
\label{sec:eshelby}
% -----------------------------------------------------------------------------

The static contribution comes from the Kelvin Green's
tensor~\eqref{eq:kelvin} via the principal-value volume
integral. Following \citet{Eshelby1957} and \citet{Mura1987},
this integral is converted to a surface integral over the
cube faces using the divergence theorem.

\subsubsection{Geometric shape constants}

On each face of the cube (e.g.\ $x_3 = a$, with $u = x_1/a$,
$v = x_2/a$), the surface integrals reduce to dimensionless
constants:
\begin{align}
  j_1 &= \int_{-1}^{1}\!\!\int_{-1}^{1}
    \frac{\dd u\,\dd v}{(1 + u^2 + v^2)^{3/2}}
  = \frac{2\pi}{3}, \label{eq:j1} \\[6pt]
  j_2 &= \int_{-1}^{1}\!\!\int_{-1}^{1}
    \frac{u^2\,\dd u\,\dd v}{(1 + u^2 + v^2)^{5/2}}
  = -\frac{2(\sqrt{3} - \pi)}{9}, \label{eq:j2} \\[6pt]
  k_1 &= \int_{-1}^{1}\!\!\int_{-1}^{1}
    \frac{\dd u\,\dd v}{(1 + u^2 + v^2)^{5/2}}
  = \frac{2(2\sqrt{3} + \pi)}{9}. \label{eq:k1}
\end{align}
These satisfy the trace constraint
\begin{equation}
  \boxed{k_1 + 2\,j_2 = \frac{2\pi}{3},}
  \label{eq:trace}
\end{equation}
which ensures the correct isotropic trace of the depolarization
tensor.

The \emph{cubic anisotropy factor}
\begin{equation}
  3\,j_2 - k_1 = -\frac{2(5\sqrt{3} - 2\pi)}{9}
  \neq 0,
  \label{eq:aniso_factor}
\end{equation}
vanishes for a sphere but is nonzero for the cube, confirming the
geometric origin of the cubic anisotropy.

\subsubsection{Static \texorpdfstring{$A$, $B$, $C$}{A, B, C}
  coefficients}

From the surface integral reduction, the static Eshelby
coefficients are:
\begin{align}
  A_\text{stat} &= -2(a_0\,j_1 + 3\,b_0\,j_2)
  = \frac{(\sqrt{3} - 2\pi)\,\alpha^2 - \sqrt{3}\,\beta^2}
         {6\pi\rho\alpha^2\beta^2},
  \label{eq:Astat} \\[6pt]
  B_\text{stat} &= 2\,b_0\,(j_1 - 3\,j_2)
  = \frac{\alpha^2 - \beta^2}
         {2\sqrt{3}\,\pi\rho\alpha^2\beta^2},
  \label{eq:Bstat} \\[6pt]
  C_\text{stat} &= 6\,b_0\,(3\,j_2 - k_1)
  = \frac{(-5\sqrt{3} + 2\pi)(\alpha^2 - \beta^2)}
         {6\pi\rho\alpha^2\beta^2}.
  \label{eq:Cstat}
\end{align}
The ratio
\begin{equation}
  \boxed{
  \frac{C_\text{stat}}{B_\text{stat}}
  = -5 + \frac{2\pi}{\sqrt{3}}
  \approx -1.372,
  }
  \label{eq:CoverB}
\end{equation}
is a \emph{pure geometric constant}, independent of all material
parameters. It depends only on the shape of the inclusion
(here, a cube) and quantifies the cubic anisotropy relative to the
shear depolarization.

% -----------------------------------------------------------------------------
\subsection{Smooth polynomial contribution}
\label{sec:smooth_contrib}
% -----------------------------------------------------------------------------

The smooth (radiation) contribution to the integral tensor comes
from integrating the smooth Green's tensor~\eqref{eq:Gsmooth} and
its second derivatives over the cube. Since $G^s_{ij}$ is a
polynomial in $x_1$, $x_2$, $x_3$, these integrals reduce to
monomial moments of the cube.

At leading order in $ka$:
\begin{align}
  A_\text{smooth} &\approx
    -\frac{2i a^3 \omega^3}{3\pi\beta^5\rho}
  + \order{a^5\omega^5},
  \label{eq:Asmooth} \\[6pt]
  B_\text{smooth} &\approx
    \frac{i a^3 \omega^3(\alpha^5 - \beta^5)}
         {3\pi\alpha^5\beta^5\rho}
  + \order{a^5\omega^5},
  \label{eq:Bsmooth} \\[6pt]
  C_\text{smooth} &\approx
    \frac{2i a^7 \omega^7(\beta^9 - \alpha^9)}
         {4725\pi\alpha^9\beta^9\rho}
  + \order{a^9\omega^9}.
  \label{eq:Csmooth}
\end{align}

\begin{quote}
\textbf{Key observation:} $C_\text{smooth}$ first appears at
$\order{(ka)^7}$, while $C_\text{stat}$ is independent of
frequency. In the Rayleigh limit ($ka \ll 1$), the ratio
$C_\text{smooth}/C_\text{stat} \sim (ka)^7$ is negligible. The
cubic anisotropy is therefore entirely dominated by the
\emph{static} Eshelby depolarization, and is a purely geometric
(frequency-independent) effect.
\end{quote}

% -----------------------------------------------------------------------------
\subsection{The scalar \texorpdfstring{$\Gamma_0$}{Gamma\_0}}
\label{sec:gamma0}
% -----------------------------------------------------------------------------

The density renormalisation requires the scalar volume integral
\begin{equation}
  \Gamma_0 = \int_V G_{ii}(\bx)\,\dd V
  = \Gamma_0^\text{stat} + \Gamma_0^\text{smooth},
  \label{eq:Gamma0}
\end{equation}
where the trace $G_{ii} = 3\,f(r) + g(r)$.

The static part involves the geometric constant
\begin{equation}
  g_0 = \int_{-1}^{1}\!\!\int_{-1}^{1}\!\!\int_{-1}^{1}
    \frac{\dd u\,\dd v\,\dd w}{\sqrt{u^2 + v^2 + w^2}}
  = -\frac{2}{3}\!\left(3\pi
    + 2\ln\!\left(70226 - 40545\sqrt{3}\right)\right),
  \label{eq:g0}
\end{equation}
an exact result from Mathematica's symbolic triple integration.
Then
\begin{equation}
  \Gamma_0^\text{stat}
  = a^2\!\left(a_0 + \frac{b_0}{3}\right) g_0
  = \frac{a^2\,(2\alpha^2 + \beta^2)}{12\pi\rho\alpha^2\beta^2}
    \,g_0.
  \label{eq:G0stat}
\end{equation}
The smooth correction $\Gamma_0^\text{smooth}$ contributes at
$\order{a^4\omega^2}$ and higher.

% =============================================================================
\section{T-Matrix Coefficients}
\label{sec:tmatrix}
% =============================================================================

The T-matrix couples the incident strain to the scattered
displacement through
\begin{equation}
  T_{mnlp} = \sum_{j,k}
    S_{mnjk}\,\Delta c_{jklp},
  \label{eq:Tcoupling}
\end{equation}
where $S_{mnjk} = \tfrac{1}{2}(I_{mjkn} + I_{njkm})$ is the
symmetrised integral tensor and
\begin{equation}
  \Delta c_{jklp} = \Dlam\,\delta_{jk}\,\delta_{lp}
    + \Dmu\,(\delta_{jl}\,\delta_{kp} + \delta_{jp}\,\delta_{kl})
  \label{eq:Dc}
\end{equation}
is the isotropic stiffness contrast
($\Dlam = \lambda' - \lambda$, $\Dmu = \mu' - \mu$).

By cubic symmetry, the T-matrix has three independent components:
\begin{equation}
  T_{mnlp} = T_1\,\delta_{mn}\,\delta_{lp}
    + T_2\,(\delta_{ml}\,\delta_{np} + \delta_{mp}\,\delta_{nl})
    + T_3\,E_{mnlp}.
  \label{eq:Tdecomp}
\end{equation}
Performing the contraction in~\eqref{eq:Tcoupling}:
\begin{equation}
  \boxed{
  \begin{aligned}
    T_1 &= (A + 4B + C)\,\Dlam + 2B\,\Dmu, \\[4pt]
    T_2 &= (A + B)\,\Dmu, \\[4pt]
    T_3 &= 2C\,\Dmu.
  \end{aligned}
  }
  \label{eq:T123}
\end{equation}

\noindent
\textbf{Physical interpretation:}
\begin{itemize}
  \item $T_1$ governs \emph{volumetric} (compression) scattering;
    it depends on both $\Dlam$ and $\Dmu$.
  \item $T_2$ governs \emph{shear} scattering (off-diagonal
    components); it depends only on $\Dmu$.
  \item $T_3$ is the \emph{cubic anisotropy} scattering; it
    vanishes for a sphere ($C = 0$) and depends only on~$\Dmu$.
\end{itemize}

\subsection{Born limit}

In the Born (static, $a \to 0$) limit, only the static Eshelby
coefficients survive, and the T-matrix becomes frequency-independent:
\begin{equation}
  T_3^\text{Born}
  = 2\,C_\text{stat}\,\Dmu
  = \frac{(-5\sqrt{3} + 2\pi)(\alpha^2 - \beta^2)\,\Dmu}
         {3\pi\rho\alpha^2\beta^2}.
  \label{eq:T3Born}
\end{equation}
This is the leading-order cubic anisotropy scattering strength,
proportional to the velocity anisotropy
$(\alpha^2 - \beta^2)$ and the shear modulus contrast~$\Dmu$.

% =============================================================================
\section{Self-Consistent Effective Medium}
\label{sec:scm}
% =============================================================================

For a random distribution of identical cubic inclusions at
concentration $\phi$, the self-consistent method
\citep{ZellerDederichs1973,Willis1981} determines the effective
medium by requiring the average T-matrix to vanish.

\subsection{Amplification factors}

Because $T_3 \neq 0$, the effective medium has \emph{four}
independent amplification factors (compared to three for the
sphere):
\begin{align}
  A_u &= \frac{1}{1 - \omega^2\Drho\,\Gamma_0},
  \label{eq:Au} \\[6pt]
  A_\theta &= \frac{1}{1 - 3T_1 - 2T_2 - T_3},
  \label{eq:Atheta} \\[6pt]
  A_e^\text{off} &= \frac{1}{1 - 2T_2},
  \label{eq:Aeoff} \\[6pt]
  A_e^\text{diag} &= \frac{1}{1 - 2T_2 - T_3}.
  \label{eq:Aediag}
\end{align}
The crucial distinction is that
$A_e^\text{off} \neq A_e^\text{diag}$ because $T_3 \neq 0$.
For a sphere, $T_3 = 0$ and both collapse to a single shear
amplification factor.

\subsection{Effective contrasts}

The effective (renormalised) contrasts are:
\begin{align}
  \Drho^* &= \Drho\,A_u, \label{eq:Drhost} \\
  \Dmu_\text{off}^* &= \frac{\Dmu}{1 - 2T_2},
  \label{eq:Dmuoff} \\
  \Dmu_\text{diag}^* &= \frac{\Dmu}{1 - 2T_2 - T_3}.
  \label{eq:Dmudiag}
\end{align}
The \emph{cubic splitting} of the effective shear modulus is
\begin{equation}
  \boxed{
  \Dmu_\text{diag}^* - \Dmu_\text{off}^*
  = \frac{\Dmu\,T_3}{(1 - 2T_2)(1 - 2T_2 - T_3)}.
  }
  \label{eq:splitting}
\end{equation}
This splitting is the observable signature of cubic anisotropy
in the effective medium: an aggregate of randomly oriented cubic
inclusions will exhibit different effective shear moduli for
diagonal and off-diagonal strain components.

% =============================================================================
\section{Numerical Verification}
\label{sec:numerical}
% =============================================================================

We verify all analytic formulae using seismological parameters:
\begin{equation*}
  \alpha = 5.0\;\text{km/s}, \quad
  \beta = 3.0\;\text{km/s}, \quad
  \rho = 2.5\;\text{g/cm}^3, \quad
  f = 10\;\text{Hz}, \quad
  a = 0.01\;\text{km}.
\end{equation*}
The background moduli are
$\mu = \rho\beta^2 = 22.5$\,GPa,
$\lambda = \rho(\alpha^2 - 2\beta^2) = 17.5$\,GPa,
and the dimensionless wavenumber products are
$k_S a = \omega a/\beta = 0.2094$,
$k_P a = \omega a/\alpha = 0.1257$.

The stiffness contrasts are $\Dlam = 1.25$\,GPa,
$\Dmu = 0.75$\,GPa.

\subsection{Integral tensor coefficients}

\begin{table}[ht]
\centering
\caption{Integral tensor coefficients $A$, $B$, $C$ (static + smooth).}
\label{tab:ABC}
\begin{tabular}{@{}lrr@{}}
\toprule
 & $\re$ & $\im$ \\
\midrule
$A$ & $-0.01220$ & $-8.61 \times 10^{-5}$ \\
$B$ & $+0.00261$ & $+3.97 \times 10^{-5}$ \\
$C$ & $-0.00359$ & $-1.04 \times 10^{-10}$ \\
\bottomrule
\end{tabular}
\end{table}

\noindent
The imaginary parts are smaller than the real parts by factors of
$10^{-2}$ (for $A$ and $B$) and $10^{-8}$ (for $C$), confirming
that the static Eshelby depolarization dominates. The
essentially zero imaginary part of~$C$ reflects
$C_\text{smooth} = \order{(ka)^7}$.

\subsection{T-matrix coefficients}

\begin{table}[ht]
\centering
\caption{T-matrix coefficients (static + smooth).}
\label{tab:Tmatrix}
\begin{tabular}{@{}lrr@{}}
\toprule
 & $\re$ & $\im$ \\
\midrule
$T_1$ & $-0.00275$ & $+1.51 \times 10^{-4}$ \\
$T_2$ & $-0.00719$ & $-3.48 \times 10^{-5}$ \\
$T_3$ & $-0.00538$ & $-1.57 \times 10^{-10}$ \\
\bottomrule
\end{tabular}
\end{table}

\noindent
Verification against the analytic formulae~\eqref{eq:T123}:
\begin{align*}
  T_1 &= (A + 4B + C)\,\Dlam + 2B\,\Dmu \\
  &= (-0.01220 + 4 \times 0.00261 - 0.00359) \times 1.25
     + 2 \times 0.00261 \times 0.75 = -0.00277 \;\checkmark \\
  T_2 &= (A + B)\,\Dmu
  = (-0.01220 + 0.00261) \times 0.75 = -0.00719 \;\checkmark \\
  T_3 &= 2C\,\Dmu
  = 2 \times (-0.00359) \times 0.75 = -0.00538 \;\checkmark
\end{align*}

\subsection{Born limit}

Setting $a \to 0$ (equivalently, retaining only the static
Eshelby terms), all dynamic corrections vanish and the
T-matrix coefficients become purely real and frequency-independent.
The Born-limit $T_3$ evaluates to
\begin{equation*}
  T_3^\text{Born}
  = \frac{(-5\sqrt{3} + 2\pi)(5^2 - 3^2) \times 0.75}
         {3\pi \times 2.5 \times 25 \times 9}
  = -0.00539,
\end{equation*}
in excellent agreement with the full result
$\re(T_3) = -0.00538$.

% =============================================================================
\section{Discussion}
\label{sec:discussion}
% =============================================================================

\subsection{Geometric origin of cubic anisotropy}

The central result of this analysis is that the cubic anisotropy
coefficient $C$ is dominated by the \emph{static} Eshelby
depolarization tensor, with the dynamic (radiation) correction
entering only at $\order{(ka)^7}$. The ratio
$C_\text{stat}/B_\text{stat} = -5 + 2\pi/\sqrt{3}$ is a pure
geometric constant, depending only on the shape of the inclusion.

This is in contrast to $A$ and~$B$, whose smooth corrections
enter at $\order{(ka)^3}$ and contribute meaningfully at moderate
frequencies. The physical reason is that the cubic anisotropy
probes the \emph{fourth} moment of the inclusion shape, while
$A$ and~$B$ probe lower moments. Higher-order shape moments
require higher-order terms in the Taylor expansion, hence higher
powers of $ka$.

\subsection{Separation of scales}

The clean separation
\begin{equation*}
  \underbrace{C_\text{stat}}_{\text{shape, real, exact}}
  + \underbrace{C_\text{smooth}}_{\order{(ka)^7},\;\text{imaginary}}
  \approx C_\text{stat}
\end{equation*}
means that in the Rayleigh limit, the cubic anisotropy
is \emph{frequency-independent}---a geometric property of the
cube that manifests identically at all frequencies for which the
Born approximation is valid.

\subsection{Implications for effective media}

For aggregates of cubic inclusions, the cubic splitting
$\Dmu_\text{diag}^* - \Dmu_\text{off}^*$~\eqref{eq:splitting}
introduces seismic anisotropy that is absent for spherical
scatterers. This provides a mechanism for apparent anisotropy in
heterogeneous earth media containing blocky or crystalline
inclusions, even when the host medium and inclusion material
are individually isotropic.

The four amplification factors~\eqref{eq:Au}--\eqref{eq:Aediag}
generalise the three-factor scheme of the self-consistent method
for spheres \citep{Korringa1973,Willis1981}. The additional degree
of freedom ($T_3$) may be constrained by independent measurements
of P and S velocity anisotropy in polycrystalline aggregates.

\bigskip

\bibliography{cubic_tmatrix}

\end{document}
